\section{Introduction}
\label{sec:intro}
\FigureArch

Politicians often answer the question they would have liked to be asked.
\citet{thomas-etal-2024-never} turned this observation into a classification task with a two-level taxonomy of response clarity and evasion, and SemEval-2026 Task~6 \citep{thomas-etal-2026-semeval} made it a shared task.
Subtask~2 asks for one of nine evasion types, and Subtask~1 for one of three clarity levels, which follow directly from the nine types.

The competition was won by pipelines that call a large language model (LLM) several times per item \citep{shi-etal-2026-teleai,peerachaidecho-rutherford-2026-chulanlp}.
Fine-tuned encoders, by contrast, stopped at about 0.50 macro-F1 on Subtask~2 no matter how large they were or how many were ensembled \citep{thomas-etal-2026-semeval}.
A leaderboard cannot tell us \emph{why} the encoders stop there, and the possible reasons predict different things.
If the labels are too noisy to learn, models should fail where annotators disagree.
If the decision rule is wrong for the metric, a better rule should help.
If the model lacks the knowledge needed to read an answer against its question, a larger pretrained model should help under the same protocol.
And if the meaning of the labels cannot be learnt from 3,448 examples, giving the model definitions should help.

We test these explanations with a deliberately controlled single-pass classifier (Figure~\ref{fig:arch}).
Every change is made on its own against a named control, every claim rests on ten paired seeds, our predictions were written down before each run, and nothing is chosen on the development set.
This gives four contributions.
\begin{itemize}
  \item A DeBERTa-v3-large system (\S\ref{sec:encoder}) whose two real improvements, longer training and the full journalist question, are separated from a series of documented failures: re-ranking, hierarchical models, rebalanced losses, per-class thresholds and model soups.
  \item Evidence that the remaining gap lies in the model and not in label noise or the decision rule (\S\ref{sec:analysis}).
  \item A one-change test of the knowledge explanation (\S\ref{sec:llm}). Replacing the backbone with Qwen3-8B-Base fine-tuned with LoRA \citep{qwen3,hu2022lora} improves both subtasks on all ten seeds and lifts the ten-seed system from 0.405 to 0.543 on Subtask~2 and from 0.648 to 0.746 on Subtask~1.
  \item An analysis of where that gain falls, which shows that our own registered prediction was wrong (\S\ref{sec:analysis}).
\end{itemize}
Because the test labels were never released, all comparisons are on the 308-item development set, and we report seed spreads, paired differences and bootstrap intervals throughout.
